\section{Explicit arithmetic on an arbitrary integral elliptic model}
\label{sec:elliptic}

Throughout this section fix integers $A,B$ with
$D_0=4A^3+27B^2\ne0$, and put
\[
 E:\quad y^2=x^3+Ax+B,\qquad S=1+|A|+|B|.
 \tag{4.1}\label{eq:4-model}
\]
Its projective equation is $Y^2Z=X^3+AXZ^2+BZ^3$.
The only point with $Z=0$ is $O=[0:1:0]$. Every other point has unique
affine coordinates, so the tagged rational codes of
Section~\ref{sec:foundations} cover all rational projective points.

\subsection{The operations and the inverse-translation identity}

Define $-O=O$ and $-(x,y)=(x,-y)$; make $O$ the identity.
For affine $P=(x,y),R=(u,v)$, if $x=u,y=-v$, put $P+R=O$.
Otherwise use
\[
 \lambda=
 \begin{cases}
  (v-y)/(u-x),&x\ne u,\\
  (3x^2+A)/(2y),&P=R,\ y\ne0,
 \end{cases}\qquad
 x'=\lambda^2-x-u,\quad y'=\lambda(x-x')-y.
 \tag{4.2}\label{eq:add}
\]
Set $P+R=(x',y')$. If $x=u$, the two curve equations imply
$(y-v)(y+v)=0$; decidable rational equality therefore leaves exactly the
listed branches. Every denominator is nonzero in its own branch.

Substitute the line $Y=\lambda(X-x)+y$ into the cubic. Its polynomial
has $x,u$ as roots, with multiplicity two in the tangent case, and the
sum of its three roots is $\lambda^2$. Polynomial division therefore
gives the third root $\lambda^2-x-u$. Reflecting its ordinate gives
\eqref{eq:add}, which proves closure by direct substitution. This proof
uses only rational arithmetic and monic division by linear factors.

We will need
\[
 (U+R)-R=U.
 \tag{4.3}\label{eq:inverse-translation}
\]
This identity has a short verification that does not require the full
associativity theorem. If the line through $U,R$ meets the cubic a third
time at $-(U+R)$, its reflection meets it at $-U,-R,U+R$.
Applying the same third-intersection rule to $U+R,-R$ therefore yields
$U$. When two of these intersections coincide, the multiplicity gives
the tangent branch. A vertical line gives the inverse branch; an input
$O$ gives the identity branch. Equivalently these are the finite case
splits of \eqref{eq:add}, verified by the same cubic factorization.
Replacing $R$ by $-R$ also gives $(U-R)+R=U$.
Thus the word-producing theorem below has an elementary arithmetic
proof independent of a separate formalization of associativity.

For the usual abelian-group formulation, \eqref{eq:add} is the
elliptic group operation. The general group laws may be reused in a
standard mathematical development; they are explicitly separate
object-theory obligations if one wants the coefficient-vector version
inside iPRA. The compact-word conclusion itself does not need them.
For the unconditional family in Section~\ref{sec:unconditional}, the
entire group law will instead be a verified four-element table.

\subsection{Square and cube denominators}

\begin{lemma}[Canonical affine form]\label{lem:denominators}
Every rational affine point of \eqref{eq:4-model} has a unique expression
\[
 x=a/d^2,\qquad y=b/d^3,\qquad
 a,b\in\Z,\ d>0,\quad \gcd(|a|,d)=1,
 \tag{4.4}\label{eq:denomform}
\]
where
\[
 b^2=a^3+Aa d^4+Bd^6.
 \tag{4.5}\label{eq:clearedcurve}
\]
The conversion is primitive recursive.
\end{lemma}
\begin{proof}
Write $x=a/c$ in reduced form, $c>0$. Then
\[
 x^3+Ax+B=(a^3+Aac^2+Bc^3)/c^3.
\]
The numerator is coprime to $c$, because it is congruent to $a^3$
modulo $c$ and $\gcd(a,c)=1$. It is consequently coprime to $c^3$.
Write $y=b/s$ in reduced form. Both $b^2/s^2$ and the preceding
fraction are reduced. Uniqueness of reduced fractions gives $s^2=c^3$
and $b^2=a^3+Aac^2+Bc^3$.

Since $c^2\mid s^2$, Lemma~\ref{lem:squares} gives $c\mid s$.
Put $d=s/c$. Then $d^2=c$, so $s=d^3$. This proves the equations.
The zero-ordinate branch causes no exception: reduced zero has
denominator $1$, and the argument gives $c=s=d=1$.
Uniqueness follows from uniqueness of the reduced $x$ denominator,
the positive square root, and then the ordinate.
\end{proof}

The uniqueness of reduced fractions used here is elementary: if
$a/c=b/d$ with both denominators positive and both pairs coprime,
$ad=bc$ and coprime cancellation imply $c\mid d$ and $d\mid c$.
Thus $c=d,a=b$. The denominator proof does not invoke valuations,
prime factorization, or algebraic number theory.

\subsection{Integer height and complete bounded-height enumeration}

For a reduced projective rational pair $[a:c]$, set
$H([a:c])=\max(|a|,|c|)$. For an affine elliptic point use its
$x$-coordinate; put $x(O)=[1:0]$. Thus
\[
 H(O)=1,\qquad H(P)=\max(|a|,d^2)
 \quad\hbox{in \eqref{eq:denomform}}.
 \tag{4.6}
\]
All heights are positive integers. Define
\[
 h(P)=\htwo H(P)=\min\{n\le H(P):H(P)\le2^n\}.
 \tag{4.7}\label{eq:binaryheight}
\]
The minimum is a bounded scan. The elementary inequality $N\le2^N$
for $N\ge1$ proves its success. In particular $h(O)=0$, and
$h(P)=0$ whenever $H(P)=1$. No real logarithm occurs in the definition.

\begin{lemma}[An explicit finite point box]\label{lem:ell-enum}
For an integer $M\ge1$, every affine point with $H(P)\le M$ occurs among
the triples
\[
 |a|\le M,\quad 1\le d\le M,\quad d^2\le M,\quad
 |b|\le SM^2,
 \tag{4.8}\label{eq:box}
\]
filtered by $\gcd(|a|,d)=1$ and \eqref{eq:clearedcurve}.
Adding $O$ gives a complete primitive recursive enumeration
$\mathcal L_E(M)$. Its length is at most
\[
 1+(2M+1)M(2SM^2+1).
 \tag{4.9}
\]
\end{lemma}
\begin{proof}
Lemma~\ref{lem:denominators} gives the form and the first three bounds.
The equation implies
$b^2\le(1+|A|+|B|)M^3=SM^3\le S^2M^4$.
Since all compared magnitudes are nonnegative, $|b|\le SM^2$.
The displayed box has the stated number of candidates before filtering.
Conversely every retained triple gives a rational point of the required
height. Each loop has its displayed numerical bound, and all tests are
decidable. Completeness is the preceding inequality, not a conclusion
drawn from an empirical search.
\end{proof}

For Theorem~\ref{thm:abstract} take $L_E(C)=\mathcal L_E(2^C)$.
One may implement the enumeration more efficiently by testing square
roots of the right side of \eqref{eq:clearedcurve}, rather than scanning
every $b$; the same bound proves completeness.

\subsection{A finite polynomial certificate for duplication}
\label{subsec:dup}

The tangent formula in \eqref{eq:add} gives the homogeneous map
\[
 x(2P)=[F(X,Z):G(X,Z)]\quad\text{when }x(P)=[X:Z],
 \tag{4.10}\label{eq:dupmap}
\]
where
\[
\begin{split}
 F(X,Z)&=X^4-2AX^2Z^2-8BXZ^3+A^2Z^4,\\
 G(X,Z)&=4X^3Z+4AXZ^3+4BZ^4.
\end{split}
\tag{4.11}\label{eq:FG}
\]
Indeed $x(2P)=(3x^2+A)^2/(4y^2)-2x$; substitute the curve equation
and homogenize. At $O$ the pair is $[1:0]$.
If $y=0$, the denominator vanishes and the numerator is
$(3x^2+A)^2$, which is nonzero by nonsingularity. Thus the same
projective formula handles every branch.

Construct the integer $8$ by $8$ matrix $M_{A,B}$ of the linear map
\[
 (U,V)\longmapsto UF+VG,
 \tag{4.12}
\]
where $U,V$ are binary cubics, and the output is a binary form of degree
$7$. Use the input order
$X^3,X^2Z,XZ^2,Z^3$ first for $U$, then for $V$, and output order
$X^7,X^6Z,\ldots,Z^7$. The columns are simply these monomials times
$F$ or $G$. Determinant expansion gives
\[
 R_0=\det M_{A,B}=256(4A^3+27B^2)^2\ne0.
 \tag{4.13}\label{eq:resultant}
\]
This is a fixed-size polynomial computation, not an appeal to a
nonconstructive elimination theorem.

For an alternative check of nonvanishing, put $f=X^3+AX+B$ and
$f'=3X^2+A$. Direct expansion proves
\[
 (27B-18AX)f+(6AX^2-9BX+4A^2)f'=4A^3+27B^2.
 \tag{4.14}
\]
Moreover $F(X,1)=f'(X)^2-8Xf(X)$ and $G(X,1)=4f(X)$.
Thus the dehomogenized polynomials are coprime over $\Q$; their
homogeneous versions have no common factor at infinity since
$F(X,0)=X^4$. This also makes the Sylvester matrix invertible by
finite polynomial division and rational linear algebra.

The first and last columns of $\adj M_{A,B}$ supply integral binary
cubics $U_X,V_X,U_Z,V_Z$ such that
\[
 U_XF+V_XG=R_0X^7,\qquad
 U_ZF+V_ZG=R_0Z^7.
 \tag{4.15}\label{eq:bezout7}
\]
All coefficients are computed by fixed-size determinants. If a positive
integer divides $R_0$ and all the coefficients of the four cubics, divide
both identities by it. Denote the resulting nonzero integer by $R_*$.
No sharp normalization is required; a common gcd suffices.
Let
\[
 L=\max\bigl(1,\norm{U_X}_1+\norm{V_X}_1,
                   \norm{U_Z}_1+\norm{V_Z}_1\bigr).
 \tag{4.16}\label{eq:L}
\]
Here $\norm{\cdot}_1$ is the sum of the absolute coefficients, after any
normalization. This is an explicit positive integer depending only on
$A,B$. With
$M=\max(1,2|A|,8|B|,A^2,4,4|A|,4|B|)$, a coarser completely
displayed upper bound is $L\le8\cdot7!\,M^7$.

\begin{lemma}[Duplication lower bound]\label{lem:dup-height}
For every rational point,
\[
                  H(P)^4\le L\,H(2P).
 \tag{4.17}\label{eq:dupheight}
\]
\end{lemma}
\begin{proof}
Write $x(P)=[a:c]$ with $\gcd(|a|,|c|)=1$. Set
\[
 H=\max(|a|,|c|),\qquad T=\max(|F(a,c)|,|G(a,c)|).
\]
The two evaluations cannot both vanish, because
\eqref{eq:bezout7}, with $R_*$ in place of $R_0$, would force
$a=c=0$. Put $g=\gcd(|F(a,c)|,|G(a,c)|)>0$.
The identities show that $g$ divides both $R_*a^7$ and $R_*c^7$.
Since $a^7,c^7$ are coprime, Bezout gives $g\mid R_*$.
Consequently $g\le|R_*|$.

Choose a coordinate of magnitude $H$ and use its corresponding identity.
Every cubic multiplier is at most its coefficient norm times $H^3$,
so
\[
 |R_*|H^7\le L H^3T,\qquad |R_*|H^4\le LT.
 \tag{4.18}
\]
After reduction, $H(2P)=T/g$. Dividing the second inequality by
$g$, and using $|R_*|/g\ge1$, gives \eqref{eq:dupheight}.
All divisions used in the last statement are exact integer divisions;
one can instead multiply through by $g$ throughout.

The gcd argument controls the cancellation in projective coordinates.
Omitting it would not prove a lower height bound.
\end{proof}

\subsection{A translation upper bound with every constant specified}
\label{subsec:translation}

For a nonidentity rational point $R$, define
\[
 T_R=\max\{4S^2H(R)^2,\ H(2R),\ 1\};
 \qquad T_O=1.
 \tag{4.19}\label{eq:TR}
\]
All these quantities are directly computable.

\begin{lemma}[Translation upper bound]\label{lem:translation}
For every pair of rational points $P,R$,
\[
                  H(P-R)\le T_R H(P)^2.
 \tag{4.20}\label{eq:translation}
\]
\end{lemma}
\begin{proof}
For nonidentity $P,R$, use their canonical forms
$P=(a/d^2,b/d^3)$ and $R=(c/e^2,w/e^3)$.
If their abscissae differ, the secant formula simplifies to
\[
 x(P-R)=
 \frac{x(P)^2x(R)+x(P)x(R)^2+
       A(x(P)+x(R))+2B+2y(P)y(R)}
      {(x(P)-x(R))^2}.
 \tag{4.21}
\]
Multiplying numerator and denominator by $d^4e^4$ gives the integer pair
\[
\begin{split}
 N={}&a^2ce^2+ac^2d^2
       +A(ad^2e^4+cd^4e^2)
       +2Bd^4e^4+2bwde,\\
 D={}&(ae^2-cd^2)^2.
\end{split}
\tag{4.22}\label{eq:ND}
\]
Put $H=H(P)$ and $J=H(R)$. The curve equation gives
$b^2\le SH^3$ and $d^2\le H$, whence
$(|b|d)^2\le SH^4\le S^2H^4$, and therefore $|b|d\le SH^2$.
Likewise $|w|e\le SJ^2$.
The first two terms of $N$ each have magnitude at most $H^2J^2$;
the $A$ terms contribute at most $2|A|H^2J^2$;
the $B$ term at most $2|B|H^2J^2$;
and the last term at most $2S^2H^2J^2$.
Thus
\[
 |N|\le2(S+S^2)H^2J^2\le4S^2H^2J^2,\qquad
 |D|\le4H^2J^2.
 \tag{4.23}
\]
Reduction of the pair can only decrease its height. This proves the bound
for distinct abscissae.

If the abscissae agree, the ordinates are equal or opposite.
Then $P-R$ is $O$ or $-2R$, whose height is at most
$\max(1,H(2R))H(P)^2$. This includes a point of order two.
If $P=O$, the needed inequality is $H(R)\le T_R$;
if $R=O$, it is $H(P)\le H(P)^2$. These follow from the definitions.
All exceptional branches have now been covered.
\end{proof}

\subsection{Unconditional finite division by two}
\label{subsec:halves}

Producing the preimages of a given rational point under doubling
is much easier than finding an unknown representative of every quotient
class. We make this distinction algorithmic.

\begin{lemma}[A complete finite rational-root scan]\label{lem:ratroots}
The rational roots of any nonzero integer polynomial of fixed degree can
be listed by a primitive recursive algorithm.
\end{lemma}
\begin{proof}
Record zero exactly when the constant coefficient vanishes, then strip
the maximal power of $X$, using a scan bounded by the degree.
If a nonconstant polynomial
$c_dX^d+\cdots+c_0$ remains with $c_dc_0\ne0$, a nonzero rational
root $m/n$ in lowest terms, $n>0$, satisfies
\[
 c_dm^d+c_{d-1}m^{d-1}n+\cdots+c_0n^d=0.
 \tag{4.24}
\]
Reducing the equation modulo $m$ gives $m\mid c_0n^d$.
Since $\gcd(|m|,n^d)=1$, $|m|\mid|c_0|$.
Similarly $n\mid|c_d|$. Scan
$-|c_0|\le m\le|c_0|$ and $1\le n\le|c_d|$, test coprimality,
and test (4.24). This is complete and sound.
If stripping leaves a constant, there are no nonzero roots.
The division tests and finite scans are covered by
Lemmas~\ref{lem:gcd} and \ref{lem:scan}.
\end{proof}

\begin{theorem}[Primitive recursive rational halving]\label{thm:halves}
There is a fixed primitive recursive function $\Halves_E(T)$ that lists
exactly the rational points $Q$ with $2Q=T$. In particular rational
divisibility of a given point by two is decidable.
\end{theorem}
\begin{proof}
First take an affine target whose reduced abscissa is $u/v$, $v>0$.
By \eqref{eq:dupmap}, the abscissa $X=x(Q)$ of a preimage satisfies
\[
\begin{split}
 p_T(X)={}&vX^4-4uX^3-2AvX^2\\
          &-(8Bv+4Au)X+(A^2v-4Bu)=0.
\end{split}
\tag{4.25}\label{eq:halvingquartic}
\]
Its leading coefficient is $v>0$, so it is nonzero.
List all its rational roots by Lemma~\ref{lem:ratroots}.
For each root $x$, calculate the reduced rational
$z=x^3+Ax+B$. If $z<0$, discard it. Otherwise test the numerator
and denominator of $z$ for exact squares. If they are squares, test
both points $(x,\sqrt z)$ and $(x,-\sqrt z)$ by the full exact equation
$2Q=T$. Keep exactly the successful points, removing duplicates if desired.

Soundness follows from that final check. For completeness, a preimage of
an affine target is affine and is not of order two; its abscissa satisfies
\eqref{eq:halvingquartic}, and its ordinate is one of the two rational
square roots just tested. Thus it is retained.

For $T=O$, include $O$ and the points $(x,0)$ where $x$ is a rational
root of $x^3+Ax+B$. These are exactly the zero-denominator branches
of doubling. The cubic is monic, so the rational-root scan is again
bounded. Every step has explicit input-dependent finite bounds.
There is no search for points of unbounded height.
\end{proof}

An empty output is a genuine decision about that particular
point. It is not a decision of the universal assertion
$\Cover_E(R)$. Pointwise decidability of a quotient equivalence relation
does not by itself supply a complete list of its classes.

\section{The supplied-cover constructive elliptic Mordell--Weil theorem}
\label{sec:cover-thm}

Let $R=(R_0,\ldots,R_{r-1})$ be a nonempty list of rational points of
$E$, and suppose \eqref{eq:cover}. Adding $O$ to the list is always
permitted. Compute $L$ from \eqref{eq:L} and $T_{R_i}$ from
\eqref{eq:TR}, then put
\[
 T=\max_{i<r}T_{R_i},\quad
 D=LT,\quad d=\htwo D,\quad C=d+2.
 \tag{5.1}\label{eq:cutoff}
\]
These constants depend only on the given equation and representative list.

\begin{lemma}[A natural-number contraction]\label{lem:bitcontract}
Whenever $P=2Q+R_i$, one has
\[
          H(Q)^4\le D H(P)^2,\qquad
          2h(Q)\le h(P)+C.
 \tag{5.2}\label{eq:bitcontract}
\]
\end{lemma}
\begin{proof}
By \eqref{eq:inverse-translation}, $2Q=P-R_i$.
Combine \eqref{eq:dupheight} and \eqref{eq:translation}.
For the second inequality, if $h(Q)=0$ there is nothing to prove.
Otherwise the definition of $h$ gives
$2^{h(Q)-1}<H(Q)$, while $D\le2^d$ and $H(P)\le2^{h(P)}$.
Comparing powers of two yields
\[
 4(h(Q)-1)<d+2h(P),
 \qquad
 4h(Q)\le d+2h(P)+3\le2h(P)+2d+4.
 \tag{5.3}
\]
Dividing this last natural-number inequality by two proves the result.
This proof can be written entirely with integer multiplication, powers and
order; no logarithm laws or real analysis are being imported.
\end{proof}

\begin{theorem}[Constructive Mordell--Weil from a supplied cover]
\label{thm:elliptic-cover}
For every $A,B\in\Z$ with $4A^3+27B^2\ne0$, and every finite valid
list $R$ satisfying $\Cover_E(R)$, the explicitly computed list
\[
                G_E(R)=R\mathbin{+\!\!+}\mathcal L_E(2^C)
 \tag{5.4}\label{eq:genlist}
\]
generates $E(\Q)$ by chord-and-tangent addition and doubling.
A fixed primitive recursive algorithm returns, for each rational point
$P$, a finite word over this list evaluating to $P$.
Its arithmetic correctness under the stated covering hypothesis has an
iPRA proof by the displayed quantifier-free invariants.

In an abelian-group presentation with the group laws proved, the same
algorithm returns nonnegative integer coefficients. With
$e=h(P)\dotminus C$ and $t=\bl(e)$, every coefficient is at most $2^t$.
The size of the list is bounded by
\[
 |R|+1+(2^{C+1}+1)2^C(2S\,2^{2C}+1).
 \tag{5.5}
\]
\end{theorem}
\begin{proof}
For an input $P$, scan $i=0,\ldots,r-1$, calculate $P-R_i$ and
$\Halves_E(P-R_i)$, and select the first point found.
This is a fixed primitive recursive operation, totalized with a failure
flag if every list is empty. By $\Cover_E(R)$ there exists a correct pair
$i,Q$; Theorem~\ref{thm:halves} proves the corresponding finite list is
nonempty. Thus the failure branch is never taken on a valid input under
the covering hypothesis. Equation \eqref{eq:inverse-translation}, in
its reverse form, gives $P=2Q+R_i$.
Lemma~\ref{lem:bitcontract} proves the contraction for whichever witness
the scan selects.

Lemma~\ref{lem:ell-enum} supplies the entire base region.
Apply Theorem~\ref{thm:abstract}, or its operation-only word version.
The bound (5.5) is (4.9) with $M=2^C$, plus the representative list.

The only unbounded existential elimination in constructing the selector's
proof comes from the hypothesis $\Cover_E(R)$ at the current point.
It proves the success of a specified finite program; it does not
define a choice function by arbitrary witness selection. All induction
formulas are the Boolean scan invariant or the forward state invariant of
Section~\ref{sec:descent}.
\end{proof}

\subsection{An alternative selector using only bounded-height enumeration}

Even without the rational-root optimization, \eqref{eq:bitcontract}
implies $h(Q)\le h(P)+C$. One may scan all $i<r$ and all
$Q\in\mathcal L_E(2^{h(P)+C})$, testing $2Q+R_i=P$.
Completeness of the enumeration and the same covering hypothesis prove
success. This route is larger computationally but shows directly that
the height bounds alone remove the need for a separately supplied selector.
The native Lean finite-search module formalizes the selection mechanism.

\subsection{A small explicit duplication certificate}

For $E_1:y^2=x^3-x$,
$F=(X^2+Z^2)^2$, $G=4XZ(X^2-Z^2)$. The following two identities
can be checked by expansion:
\[
\begin{split}
 (16X^3-12XZ^2)F+(-5X^2Z-3Z^3)G&=16X^7,\\
 (-12X^2Z+16Z^3)F+(3X^3+5XZ^2)G&=16Z^7.
\end{split}
\tag{5.6}
\]
They give $L=36$. With the four points later proved to be all of $E_1(\Q)$,
$T=16$ and $C=\htwo(576)+2=12$. This is a deliberately loose general
bound: the unconditional proof below gives the exact height bound $H=1$.
The example illustrates that the general constants are explicit, rather
than asserting that their large enumeration is a practical way to solve
this curve.

\subsection{The abstract version for number fields and abelian varieties}

For a fixed number field one may encode $K=\Q[T]/(f)$ by coefficient
vectors of degree less than $\deg f$, with $f$ a supplied irreducible
polynomial. Reduction and field operations then have finite algebraic
implementations. A projective variety can be represented by finitely many
coordinate charts and polynomial membership tests, with fixed formulas
and guards for its group operations. To apply
Theorem~\ref{thm:abstract} to $A(K)$ one must additionally provide:
a natural-valued primitive recursive height; complete finite
bounded-height enumerations; an actual covering list modulo multiplication;
and a primitive recursive decomposition with the proved contraction.

The same arithmetic theorem then produces generators. This is a
relative interface, not a completed construction of heights,
number-field geometry, weak descent, or representatives for arbitrary
$A,K$. Merely supplying defining equations for an abelian variety does
not discharge those obligations. Section~\ref{sec:unconditional} supplies
actual abelian varieties of every finite dimension for which all the
required rational-point data are settled without those extra hypotheses.
